Now, if @@M0@@ is a @@M1@@-ring, the diagram:
\sgaStableTarget{sga6.e05.eq.7-5-4}@@M2@@
is commutative by 7.3, since @@M3@@ is a @@M4@@-homomorphism. The desired relations follow at once.

ii) If @@M5@@ is of characteristic zero, the ``logarithmic derivative'' map is an isomorphism from @@M6@@ onto @@M7@@, since one can integrate. Consequently, two elements of @@M8@@ are equal if and only if, for every @@M9@@, they have the same image under @@M10@@.

To show that @@M11@@ is a ring homomorphism, we must show the commutativity of the left square of diagram (7.5.2), and by the preceding remark, it suffices to show that for every @@M12@@, diagram (7.5.2) is commutative, since the right square always is. Now these are precisely the relations @@M13@@.

To show that @@M14@@ is a @@M15@@-homomorphism, it suffices to show that diagram (7.5.4) is commutative for every @@M16@@, by 7.3. Now, by the preceding remark, this is equivalent to showing for all @@M17@@ that @@M18@@, that is, the relations @@M19@@.

Finally, the relation @@M20@@ is obtained by integration from the relations @@M21@@.

\sgaStableTarget{sga6.e05.bibliography}\subsection*{BIBLIOGRAPHY}
\begin{enumerate}[label={[\arabic*]}]
\item A. Grothendieck, La théorie des classes de Chern, Bull. Soc. math. France, 86, 1958, p. 137 à 154.
\item F. Hirzebruch, Neue topologische Methoden in der algebraischen Geometrie.
\end{enumerate}

% SGA_NATIVE_CONTINUATION_364
\providecommand{\Sym}{\operatorname{Sym}}
\providecommand{\Proj}{\operatorname{Proj}}
\providecommand{\Projz}{\operatorname{Proj}_0}
\providecommand{\Spec}{\operatorname{Spec}}
\providecommand{\Comp}{\operatorname{Comp}}
\providecommand{\Dplus}{D_{+}}
\providecommand{\id}{\operatorname{id}}
\providecommand{\cG}{\mathcal G}
\providecommand{\cA}{\mathcal A}
\providecommand{\cB}{\mathcal B}
\providecommand{\cL}{\mathscr L}
\providecommand{\cM}{\mathscr M}
\providecommand{\cN}{\mathscr N}
\providecommand{\cE}{\mathscr E}
\providecommand{\cK}{\mathscr K}
\providecommand{\cP}{\mathscr P}
\providecommand{\cQ}{\mathscr Q}
\providecommand{\OO}{\mathcal O}
\providecommand{\rk}{\operatorname{rk}}
\providecommand{\colim}{\operatorname*{colim}}

% SGA_REVIEWED_FRENCH_CONTINUATION_364
% Exposé VI projective-bundle macro additions
\providecommand{\Sym}{\operatorname{Sym}}
\providecommand{\Proj}{\operatorname{Proj}}
\providecommand{\Projz}{\operatorname{Proj}_0}
\providecommand{\Spec}{\operatorname{Spec}}
\providecommand{\Comp}{\operatorname{Comp}}
\providecommand{\Dplus}{D_{+}}
\renewcommand{\id}{\operatorname{id}}
\providecommand{\cG}{\mathcal G}
\providecommand{\cA}{\mathcal A}
\providecommand{\cB}{\mathcal B}
\providecommand{\cL}{\mathscr L}
\providecommand{\cM}{\mathscr M}
\providecommand{\cN}{\mathscr N}
\providecommand{\cE}{\mathscr E}
\providecommand{\cK}{\mathscr K}
\providecommand{\cP}{\mathscr P}
\providecommand{\cQ}{\mathscr Q}
\providecommand{\OO}{\mathcal O}
\providecommand{\rk}{\operatorname{rk}}
\providecommand{\colim}{\operatorname*{colim}}

